\chapter{Vanilla Rates Options}\label{ch:rc:vanilla-rates-options}

In June 2014 the ECB set its deposit rate at $-0.10\%$, and within two years
short euro rates and many forward rates were below zero. Black's model, in
which a rate is lognormal, cannot price an option on a negative forward, or
struck at a negative rate: its formula takes the logarithm of both. For more
than twenty years caps and swaptions had been quoted as Black volatilities; the
quotes that could not exist began to appear as normal volatilities, in basis
points a year, or as volatilities of a shifted rate. The prices did not care
which model was used to quote them; the risk numbers did. This chapter prices
the two vanilla rates options, caplets and swaptions, in both models, strips
the volatility of each caplet out of the quoted caps, and shows what the
settlement terms of a swaption do to its price. The products and their
market are those of One Quant Book 2, chapter 13; the curves are chapter~2's
euro curves.

\section{Black and Bachelier for rates}

A caplet on six-month Euribor fixing at $T_{k-1}$ and paid at $T_k$ pays
$N\delta_k\max(F_k(T_{k-1})-K,0)$, where $F_k(t)$ is the forward of that fixing.
Under the $T_k$-forward measure $\mathbb Q^{T_k}$ of the \omterm{def:rc:multi-curve-and-collateral-discounting:curves}{discount curve} (One
Quant Book 4, chapter 5), $F_k$ is a martingale, and so is the par rate
$S_{a,b}$ of a swap under the annuity measure $\mathbb Q^{a,b}$. Both options are
therefore options on a martingale, and the model only has to say how that
martingale is distributed.

\begin{proposition}[Caplets and swaptions under their natural measures]\label{prop:rc:vanilla-rates-options:measures}
\[
\mathrm{Caplet}_0 = N\delta_kP(0,T_k)\,\E^{T_k}\bigl[(F_k(T_{k-1})-K)^+\bigr],
\qquad
\mathrm{Payer}_0 = N A_{a,b}(0)\,\E^{a,b}\bigl[(S_{a,b}(T_a)-K)^+\bigr].
\]
If $dF = \sigma_N\,dW$ the expectation is Bachelier's formula (One Quant Book 2,
chapter 13); if $dF = \sigma F\,dW$ it is Black's formula (One Quant Book 5,
chapter 3), with the forward in place of the spot and no drift.
\end{proposition}
\begin{proof}[Partial proof]
Take $P(t,T_k)$ as numeraire: the caplet's price divided by it is a martingale,
and at $T_k$ the numeraire is one. $F_k(t) = (P(t,T_{k-1})/P(t,T_k)-1)/\delta_k$
is a traded asset divided by the numeraire, hence a martingale. The same
argument with the annuity as numeraire gives the swaption; with multiple curves,
the projected forward is defined as the $\mathbb Q^{T_k}$ expectation of the
fixing.
\end{proof}

\begin{definition}[Lognormal volatility]\label{def:rc:vanilla-rates-options:lognormal}
The \emph{lognormal volatility}\index{lognormal volatility} $\sigma_B$ of a
rates option is the volatility at which Black's formula, applied to the
forward rate, reproduces its price; it is a relative volatility, a percentage
of the rate, and exists only for positive forward and strike.
\end{definition}

\begin{proposition}[Normal and lognormal volatilities at the money]\label{prop:rc:vanilla-rates-options:atm}
For an option at the money ($K=F>0$), the two models give the same price when
\[
\sigma_N\sqrt{T}\,\frac{1}{\sqrt{2\pi}} = F\bigl(2\Phi(\tfrac12\sigma_B\sqrt T)-1\bigr),
\quad\text{so}\quad
\sigma_N \approx \sigma_B F \ \text{ for small }\sigma_B\sqrt T.
\]
\end{proposition}
\begin{proof}
The two at-the-money prices are $\sigma_N\sqrt{T}/\sqrt{2\pi}$ and
$F(2\Phi(\sigma_B\sqrt T/2)-1)$; expand $\Phi$ to first order.
\end{proof}

A normal model priced at one volatility for all strikes is a steep negative
skew in lognormal terms (\cref{fig:rc:vanilla-rates-options:smile}): for a
five-year option on a forward of 2.50\% at 80 basis points of normal
volatility, the Black volatility is 32.7\% at the money (the approximation
gives $80/250 = 32\%$), 71.1\% at a strike of 0.50\% and 23.8\% at 4.50\%. A
market whose smile is flat in normal terms therefore looks strongly skewed to
anyone quoting in Black terms, and a model choice is also a choice of how the
smile moves when rates move.

\begin{omfigure}
\begin{tikzpicture}
\begin{axis}[width=11cm,height=5.4cm,
  xlabel={strike (\%)},ylabel={lognormal volatility (\%)},
  tick label style={font=\small},label style={font=\small},
  xmin=0.5,xmax=5,ymin=15,ymax=75,grid=major,grid style={gray!25},
  table/col sep=comma]
\addplot[omDef,very thick] table[x=k,y=black]{figdata/rates-credit-risk/04-vanilla-rates-options/smile.csv};
\draw[gray,dashed] (axis cs:2.5,15) -- (axis cs:2.5,75) node[pos=0.9,right,font=\footnotesize,text=black] {forward 2.50\%};
\end{axis}
\end{tikzpicture}

\omcaption{The lognormal (Black) implied volatility of five-year options on a
forward of 2.50\% priced with a single normal volatility of 80 basis points.
A flat normal smile is a steep lognormal skew. Data: the chapter's
tutorial.}\label{fig:rc:vanilla-rates-options:smile}
\end{omfigure}

\begin{dated}{2026-09}{Negative rates and the conventions of rates options}\label{dat:rc:vanilla-rates-options:negative}
The ECB's deposit facility rate was negative from 11 June 2014 ($-0.10\%$) to 27
July 2022, with a low of $-0.50\%$ from 18 September 2019. Options on negative
forwards or strikes have no Black volatility, and the market began quoting caps,
floors and swaptions also in normal or in shifted Black volatilities
(chapter~5). Since 26 November 2018 the default cash-settlement method for
euro swaptions in the ISDA settlement matrix has been the collateralised cash
price instead of the par-yield method.
\end{dated}

\section{Caps, floors and caplet stripping}

A cap is quoted as one number, its price or a single volatility for all its
caplets. The caplets of a five-year cap are the same caplets as the first four
years of a ten-year cap, and cannot have two prices.

\begin{definition}[Flat volatility, caplet volatility]\label{def:rc:vanilla-rates-options:flat}
The \emph{flat volatility}\index{flat volatility} of a cap is the single
volatility at which the sum of its caplets' prices equals the cap's price. A
\emph{caplet volatility}\index{caplet volatility} is the volatility of one
caplet, a function of its fixing date (and strike); the caplet volatilities of
all caps of one strike form a single term structure.
\end{definition}

\begin{definition}[Caplet stripping]\label{def:rc:vanilla-rates-options:strip}
\emph{Caplet stripping}\index{caplet stripping} is the calibration of a term
structure of caplet volatilities that reprices every quoted cap of a strike:
with caps of increasing maturity, the caplets between two consecutive
maturities are given one volatility, solved so that the longer cap reprices,
the shorter caplets' volatilities held.
\end{definition}

\omcode{code/firm/capfloor/firm_capfloor.py}{63}{86}{Stripping piecewise-constant caplet volatilities from flat cap volatilities.}\label{lst:rc:vanilla-rates-options:strip}

\begin{example}[A euro cap strip]\label{ex:rc:vanilla-rates-options:strip}
Flat normal volatilities of caps struck at 3\% on six-month Euribor of 58, 66,
72, 75, 76, 75 and 72 basis points at one, two, three, four, five, seven and
ten years (illustrative; the first caplet, already fixed, excluded) strip into
caplet volatilities of 58.0, 66.6, 75.7, 78.9, 78.0, 73.7 and 67.9 basis points
on the successive intervals (\cref{fig:rc:vanilla-rates-options:strip}). Where
flat volatilities rise, the new caplets must be above the flat level to lift the
average; the hump of the caplet curve is higher and earlier than that of the
flat curve.
\end{example}

\begin{omfigure}
\begin{tikzpicture}
\begin{axis}[width=12cm,height=5.4cm,
  xlabel={caplet fixing time (years)},ylabel={normal volatility (bp)},
  tick label style={font=\small},label style={font=\small},
  xmin=0,xmax=10,ymin=50,ymax=85,grid=major,grid style={gray!25},
  legend columns=2,legend style={font=\footnotesize,at={(0.5,-0.25)},anchor=north,draw=none,fill=none,/tikz/every even column/.append style={column sep=8pt}},
  table/col sep=comma]
\addplot[omThm,very thick,const plot,mark=*,mark size=1.3pt] table[x=t,y=caplet]{figdata/rates-credit-risk/04-vanilla-rates-options/strip.csv};
\addplot[omDef,thick,dashed,const plot] table[x=t,y=flat]{figdata/rates-credit-risk/04-vanilla-rates-options/strip.csv};
\legend{stripped caplet volatility,flat volatility of the shortest cap containing the caplet}
\end{axis}
\end{tikzpicture}

\omcaption{Caplet volatilities stripped from the flat volatilities of caps
struck at 3\%, one per interval between quoted maturities. Data: the chapter's
illustrative quotes on chapter~2's euro curves; the chapter's
tutorial.}\label{fig:rc:vanilla-rates-options:strip}
\end{omfigure}

\begin{omfigure}
\begin{tikzpicture}[x=1.15cm,y=1cm,font=\small]
\draw[->,thick] (0,0) -- (10.6,0) node[right] {time};
\foreach \x/\l in {0/spot,1/$T_0$,2/$T_1$,3/$T_2$,4/$T_3$,9/$T_{n-1}$,10/$T_n$} {\draw (\x,-0.1) -- (\x,0.1); \node[below,font=\footnotesize] at (\x,-0.12) {\l};}
\node at (6.5,0.35) {$\dots$};
\foreach \a/\b/\y in {1/2/1.1,2/3/1.1,3/4/1.1,9/10/1.1} {\draw[omDef,thick,fill=omDef!12] (\a,\y) rectangle (\b,\y+0.45);}
\node[font=\footnotesize] at (1.5,1.32) {1};
\node[font=\footnotesize] at (2.5,1.32) {2};
\node[font=\footnotesize] at (3.5,1.32) {3};
\node[font=\footnotesize] at (9.5,1.32) {$n$};
\node[font=\footnotesize,anchor=west] at (4.3,1.32) {caplets};
\draw[dashed,gray,thick] (0,1.1) rectangle (1,1.55); \node[font=\footnotesize,gray] at (0.5,1.32) {fixed};
\draw[-{Stealth[length=2mm]},omThm] (2,1.05) -- (2,0.15);
\draw[-{Stealth[length=2mm]},omExo] (3,1.05) -- (3,0.15);
\node[font=\footnotesize,text=omThm,anchor=east] at (1.8,2.3) {caplet 2 fixes at $T_1$};
\node[font=\footnotesize,text=omExo,anchor=west] at (3.3,2.3) {and pays at $T_2$};
\draw[omThm,thin] (1.8,2.3) -- (2,1.6);
\draw[omExo,thin] (3.3,2.3) -- (3,1.6);
\end{tikzpicture}

\omcaption{A cap on six-month Euribor from spot: each caplet's rate fixes at
the start of its period and is paid at the end. The first period's rate is
already fixed at trade, so the market's caps start with the second caplet.}\label{fig:rc:vanilla-rates-options:timeline}
\end{omfigure}

\section{Swaptions: physical and cash settlement}

A physically settled swaption delivers the swap; its value is the annuity
times an option on the forward swap rate
(\cref{prop:rc:vanilla-rates-options:measures}). In euros and sterling most
swaptions are settled in cash, and the settlement amount needs an annuity that
both parties can compute on the exercise date from one number.

\begin{definition}[Cash-settled swaption, cash annuity]\label{def:rc:vanilla-rates-options:cash}
A \emph{cash-settled swaption}\index{cash-settled swaption} pays, on
exercise, the value of the underlying swap computed by an agreed method
instead of delivering it. Under the par-yield method the value is
$N\,a(S_T)\,(S_T-K)^+$ with the \emph{cash annuity}\index{cash annuity}
\[
a(S) = \sum_{i=1}^{nm}\frac{1/m}{(1+S/m)^i},
\]
the annuity of the swap discounted at its own par rate $S$ ($m$ payments a
year); under the collateralised-cash-price method it is the swap's value on the
\omterm{def:rc:multi-curve-and-collateral-discounting:curves}{discount curve} of the agreed clearing house, as if physically settled.
\end{definition}

\omcode{code/firm/capfloor/firm_capfloor.py}{110}{121}{The cash annuity and the traditional par-yield cash-settlement formula.}\label{lst:rc:vanilla-rates-options:cash}

\begin{example}[Five years into ten]\label{ex:rc:vanilla-rates-options:swaption}
On chapter~2's curves the ten-year swap starting in five years has a forward
par rate of 3.098\% against six-month Euribor and a (discount-curve) annuity of
7.740. At 80 basis points of normal volatility (a \omterm{def:rc:vanilla-rates-options:lognormal}{lognormal volatility} of
26.2\%), the at-the-money physical payer on EUR~100 million is worth
EUR~5\,525\,150. The traditional cash formula, $P(0,T)\,a(S_0)\times$ option,
gives EUR~5\,410\,452, 2.1\% less: $P(0,T)\,a(S_0) = 7.579$, because the \omterm{def:rc:vanilla-rates-options:cash}{cash
annuity} discounts at the Euribor swap rate, above the €STR rates of the
collateralised annuity. The par-yield payoff is also not a function of the
annuity measure's numeraire, so the formula is itself an approximation; the
collateralised cash price removes both problems by settling at the physical
value.
\end{example}

\section{Vega, the smile, and the choice of model}

The normal vega of a caplet, $\partial\mathrm{Caplet}/\partial\sigma_N =
N\delta P\sqrt T\varphi(d)$ with $d=(F-K)/(\sigma_N\sqrt T)$, is largest at the
money and does not depend on the level of rates; the Black vega is
proportional to $F$. A desk that hedges in one model and is marked in the other
sees vega move when rates move with no change in the market. Two rules follow.
Hedge in the model in which the market's smile is most stable (for euro and
dollar rates, normal or shifted models since the negative-rate years); and
report vega per basis point of normal volatility, which is additive across
caplets and swaptions of different forwards.

\begin{omfigure}
\begin{tikzpicture}
\begin{axis}[width=11cm,height=5.2cm,
  xlabel={cap strike (\%)},ylabel={premium (bp a year)},
  tick label style={font=\small},label style={font=\small},
  xmin=2,xmax=4.5,ymin=0,ymax=75,grid=major,grid style={gray!25},
  table/col sep=comma]
\addplot[omDef,very thick,mark=*,mark size=1.4pt] table[x=k,y=bp]{figdata/rates-credit-risk/04-vanilla-rates-options/capstrike.csv};
\draw[gray,dashed] (axis cs:3,0) -- (axis cs:3,75) node[pos=0.88,right,font=\footnotesize,text=black] {the treasurer's 3\%};
\end{axis}
\end{tikzpicture}

\omcaption{Running cost of a five-year cap on six-month Euribor on chapter~2's
curves, by strike, with the caplet volatilities stripped at 3\% used at every
strike (the smile ignored). The cost falls from 70 basis points a year at a
strike of 2\% to 4 at 4.5\%. Data: the chapter's tutorial.}\label{fig:rc:vanilla-rates-options:capstrike}
\end{omfigure}

\begin{remark}[Flat or stripped: same price, different risk]\label{rem:rc:vanilla-rates-options:samerisk}
The five-year cap of the problem below has the same price at its \omterm{def:rc:vanilla-rates-options:flat}{flat
volatility} and at the stripped caplet volatilities, by construction; its total
vega is the same too, but its bucketed vega is not: the stripped curve puts it
on the caplets' own expiries, where the hedging caps and swaptions sit.
\end{remark}

\section{Tutorial: strip, price, compare}\label{tut:rc:vanilla-rates-options:tutorial}

\begin{tutorial}
\textbf{Goal.} Strip caplet volatilities from cap quotes, price a cap and a
swaption on chapter~2's euro curves, and compare the settlement methods.
\textbf{End state:}
\cref{fig:rc:vanilla-rates-options:smile,fig:rc:vanilla-rates-options:strip} and the
numbers of \cref{ex:rc:vanilla-rates-options:swaption}.
\begin{tutsteps}
\item \textbf{Smile}: \texttt{smile\_from\_flat\_normal()} converts a flat normal
volatility into Black volatilities by strike with Book~2's
\texttt{normal\_to\_black}.
\item \textbf{Strip}: \texttt{stripped()} returns seven knots; check that every
cap reprices at its \omterm{def:rc:vanilla-rates-options:flat}{flat volatility}.
\item \textbf{Swaption}: \texttt{swaption\_table()} gives the forward, the two
annuities and the two prices.
\item \textbf{Cap}: \texttt{treasurer\_cap()} and \texttt{cap\_strike\_table()};
\texttt{fig\_rc\_vanilla.py} writes the charts.
\end{tutsteps}
\textbf{What to change next.} Price the cap with a flat \omterm{def:rc:vanilla-rates-options:lognormal}{lognormal volatility} of
25\% and compare; strip with caps quoted every year instead of at one, two,
three, four, five, seven and ten years.
\end{tutorial}

\section{Build: caps, floors and swaptions}\label{bld:rc:vanilla-rates-options:capfloor}

\begin{build}
\bfield{Purpose} The vanilla rates options of the miniature firm, priced on the
multi-curve; the calibration instruments of the models of chapters 5 to 10, and
the options the risk engine of \cref{ch:rc:build-a-risk-engine} registers with
the pricing library.
\bfield{Interface} \texttt{Cap(start, years, strike, months, floor, skip\_first)};
\texttt{caplets}, \texttt{cap\_price(cap, proj, disc, vol, model)},
\texttt{cap\_vega}, \texttt{strip\_caplet\_vols}, \texttt{piecewise};
\texttt{forward\_swap}, \texttt{swaption\_physical}, \texttt{cash\_annuity},
\texttt{swaption\_cash\_irr}.
\bfield{Rules} Forward-looking caplets fixed at the start of their period; the
first caplet skipped by default; \texttt{model} is \texttt{normal} or
\texttt{lognormal}; Book 2's \texttt{firm\_normalvol} imported, not edited.
\bfield{Acceptance tests} \texttt{code/firm/capfloor/tests/}: zero volatility
gives intrinsic value; cap minus floor is the forward swap; stripped
volatilities reprice every cap; payer minus receiver is the forward swap; the
\omterm{def:rc:vanilla-rates-options:cash}{cash annuity} has the closed form and falls with the rate.
\bfield{Stretch} Strike-dependent stripping (a caplet smile per expiry); the
collateralised cash price with a separate settlement curve; backward-looking
caplets (chapter~10).
\end{build}

\begin{omsources}
\item F.\ Black, ``The pricing of commodity contracts'', \emph{Journal of
Financial Economics} 3, 1976.
\item ISDA, ``Market practice change for settlement of EUR swaptions to
collateralized cash price'', November 2018.
\item D.\ Brigo and F.\ Mercurio, \emph{Interest Rate Models: Theory and
Practice}, 2nd ed., Springer, 2006, chapter 6.
\item ECB, \emph{Key ECB interest rates}.
\end{omsources}

\section{Exercises}

\begin{exercise}[$\star$]\label{exo:rc:vanilla-rates-options:1}
A caplet on EUR~10 million struck at 3\% fixes at 3.40\% for a period with
accrual 0.5069. What does it pay, and when?
\end{exercise}

\begin{exercise}[$\star$]\label{exo:rc:vanilla-rates-options:2}
A five-year at-the-money option on a forward of 2.50\% has a normal volatility
of 80 basis points. Estimate its \omterm{def:rc:vanilla-rates-options:lognormal}{lognormal volatility}, and compare with the exact
32.7\%.
\end{exercise}

\begin{exercise}[$\star$]\label{exo:rc:vanilla-rates-options:3}
Why do market caps on six-month Euribor exclude the first caplet?
\end{exercise}

\begin{exercise}[$\star\star$]\label{exo:rc:vanilla-rates-options:4}
The one-year cap has a \omterm{def:rc:vanilla-rates-options:flat}{flat volatility} of 58 and the two-year of 66 basis
points. Without computing, is the stripped volatility of the second-year caplets
above or below 66? Check against \cref{ex:rc:vanilla-rates-options:strip}.
\end{exercise}

\begin{exercise}[$\star\star$]\label{exo:rc:vanilla-rates-options:5}
Explain the 2.1\% gap between the par-yield cash formula and the physical price
in \cref{ex:rc:vanilla-rates-options:swaption}, and why the collateralised cash
price closes it.
\end{exercise}

\begin{exercise}[$\star\star$]\label{exo:rc:vanilla-rates-options:6}
Read \cref{fig:rc:vanilla-rates-options:smile}: why does a flat normal smile
look like a negative lognormal skew? What smile would a flat lognormal market
look like in normal terms?
\end{exercise}

\begin{exercise}[$\star\star\star$]\label{exo:rc:vanilla-rates-options:7}
\emph{Coding.} Find the floor strike that makes a five-year collar (long the 3\%
cap, short the floor) cost nothing, with the stripped volatilities at every
strike.
\end{exercise}

\begin{exercise}[$\star\star\star$]\label{exo:rc:vanilla-rates-options:8}
\emph{Find the flaw.} ``The five-year cap is quoted at 76 basis points, so I
price the two caplets fixing after four years at 76.'' What is wrong, and by
how much does it misprice them on EUR~200 million?
\end{exercise}

\section{Problem: The Treasurer's Cap}

\begin{problem}[{Weekend problem --- capping a floating-rate loan}]\label{pb:rc:vanilla-rates-options:1}
A company has borrowed EUR~200 million for five years at six-month Euribor plus
a margin, first rate already fixed. Its board wants the rate capped at 3\%. On
chapter~2's curves and the stripped volatilities of
\cref{ex:rc:vanilla-rates-options:strip} it buys the five-year cap at 3\%.

\medskip\noindent\textbf{Part I --- The cap.}
\begin{enumerate}
\item How many caplets does it hold, and why that number?
\item Give its premium in euros and as a percentage of notional.
\item Express the premium as a running cost in basis points a year.
\item Price it at the five-year \omterm{def:rc:vanilla-rates-options:flat}{flat volatility} instead, and explain the result.
\item The average forward of the caplets is 2.44\%. Why pay anything for a cap
struck above it?
\end{enumerate}

\medskip\noindent\textbf{Part II --- Risk.}
\begin{enumerate}[resume]
\item Give its vega per basis point of normal volatility.
\item If normal volatilities rise by 10 basis points, how much does the cap gain?
\item Would its lognormal vega grow or shrink if rates rose, at constant normal
volatility?
\item Who is short this vega, and how is it hedged?
\item Price the cap at a flat \omterm{def:rc:vanilla-rates-options:lognormal}{lognormal volatility} of 25\% and comment.
\end{enumerate}

\medskip\noindent\textbf{Part III --- Scenarios.}
\begin{enumerate}[resume]
\item If every remaining fixing is 3.50\%, what does the cap pay in total,
undiscounted?
\item Give the present value of those payments today.
\item At what average fixing does the company break even on the premium,
approximately?
\item If all fixings stay below 3\%, what has the company bought?
\item Why is the cap not an option on the average rate?
\end{enumerate}

\medskip\noindent\textbf{Part IV --- Judgement.}
\begin{enumerate}[resume]
\item Compare the cap with paying fixed on a five-year swap.
\item Give the floor strike of a zero-cost collar and say what the company gives
up.
\item Why might a bank quote this cap above the model price?
\item State the \emph{named result}: the premium, its running cost and the vega
of the treasurer's cap.
\item In one sentence: what does a cap buy that a swap does not?
\end{enumerate}
\end{problem}

\section{Interview questions}

\begin{interviewq}[$\star$ \iqroles{trader, bank}]\label{iq:rc:vanilla-rates-options:1}
What is a caplet, and under which measure is its forward rate a martingale?
\end{interviewq}

\begin{interviewq}[$\star\star$ \iqroles{researcher}]\label{iq:rc:vanilla-rates-options:2}
Derive the relation between normal and lognormal at-the-money volatility.
\end{interviewq}

\begin{interviewq}[$\star\star$ \iqroles{researcher, developer}]\label{iq:rc:vanilla-rates-options:3}
How do you strip caplet volatilities from cap quotes, and what can go wrong?
\end{interviewq}

\begin{interviewq}[$\star\star$ \iqroles{trader}]\label{iq:rc:vanilla-rates-options:4}
Why would a desk prefer to risk-manage in normal volatility?
\end{interviewq}

\begin{interviewq}[$\star\star\star$ \iqroles{researcher, bank}]\label{iq:rc:vanilla-rates-options:5}
What is the problem with pricing a par-yield \omterm{def:rc:vanilla-rates-options:cash}{cash-settled swaption} with Black's
formula and the \omterm{def:rc:vanilla-rates-options:cash}{cash annuity}, and how did the market fix it?
\end{interviewq}

\begin{interviewq}[$\star\star\star$ \iqroles{developer, risk}]\label{iq:rc:vanilla-rates-options:6}
Your caplet pricer returns NaN for some trades after rates fell. Why, and what
do you change?
\end{interviewq}
